\section{Remember Me at the Altar of the Lord}\label{sec:suffrage}

The Church prayed for the dead from the beginning, and her teaching on purgatory rests in great part on that prayer. The Catechism says that the teaching on purification ``is also based on the practice of prayer for the dead, already mentioned in Sacred Scripture,'' and it points to Judas Machabeus, who after the battle sent an offering to Jerusalem ``for sacrifice to be offered for the sins of the dead'' (2~Mach 12:43).\endnote{\emph{Catechism of the Catholic Church} 1032, official English.} The sacred writer comments: ``It is therefore a holy and wholesome thought to pray for the dead, that they may be loosed from sins'' (12:46).

\subsection{The witness of the Fathers}

The earliest Fathers speak of the practice as something every Christian already knows. Tertullian lists it among the customs that the Church keeps by tradition, beside the renunciations of baptism and the sign of the cross: ``As often as the anniversary comes round, we make offerings for the dead as birthday honours.''\endnote{Tertullian, \emph{De corona} 3, trans.\ S. Thelwall, \emph{Ante-Nicene Fathers}, vol.~3, New Advent edition, \url{https://www.newadvent.org/fathers/0304.htm} (accessed 16 October 2026).} Year by year, on the anniversary of a Christian's death, the Church offered for him.

Cyril of Jerusalem, explaining the Eucharistic liturgy to the newly baptized, describes the commemoration of the departed that follows, ``after the spiritual sacrifice, the bloodless service, is completed'':

\begin{quote}
Then we commemorate also those who have fallen asleep before us, first Patriarchs, Prophets, Apostles, Martyrs, that at their prayers and intercessions God would receive our petition. Then on behalf also of the Holy Fathers and Bishops who have fallen asleep before us, and in a word of all who in past years have fallen asleep among us, believing that it will be a very great benefit to the souls, for whom the supplication is put up, while that holy and most awful sacrifice is set forth.\endnote{Cyril of Jerusalem, \emph{Mystagogical Catecheses} 5.8--9, trans.\ E.~H. Gifford, \emph{Nicene and Post-Nicene Fathers}, second series, vol.~7.}
\end{quote}

To those who asked what such a commemoration could profit a soul that had died in sin, Cyril answered with a parable. If a king had banished certain men who had offended him, and their kinsmen wove a crown and offered it to him on their behalf, would he not grant a remission of their penalties? ``In the same way we, when we offer to Him our supplications for those who have fallen asleep, though they be sinners, weave no crown, but offer up Christ sacrificed for our sins, propitiating our merciful God for them as well as for ourselves.''\endnote{Cyril, \emph{Mystagogical Catecheses} 5.10.}

John Chrysostom appeals to the same rite as an apostolic ordinance: ``Not in vain did the Apostles order that remembrance should be made of the dead in the dreadful Mysteries. They know that great gain resulteth to them, great benefit.''\endnote{John Chrysostom, \emph{Homilies on Philippians}, Homily~3, Oxford translation revised by John A. Broadus, \emph{Nicene and Post-Nicene Fathers}, first series, vol.~13.} And to mourners who asked what they could still do for their dead, he gave an answer that the Catechism quotes:

\begin{quote}
Let us then give them aid and perform commemoration for them. For if the children of Job were purged by the sacrifice of their father, why dost thou doubt that when we too offer for the departed, some consolation arises to them? \ldots\ Let us not then be weary in giving aid to the departed, both by offering on their behalf and obtaining prayers for them.\endnote{John Chrysostom, \emph{Homilies on First Corinthians}, Homily~41, 8, trans.\ Talbot W. Chambers, \emph{Nicene and Post-Nicene Fathers}, first series, vol.~12; quoted at \emph{Catechism of the Catholic Church} 1032.}
\end{quote}

\noindent ``Why therefore dost thou grieve?'' he asks. ``Why mourn, when it is in thy power to gather so much pardon for the departed?''

Augustine states whom these offerings help. ``Nor can it be denied that the souls of the dead are benefited by the piety of their living friends, who offer the sacrifice of the Mediator, or give alms in the church on their behalf.'' But these services help only those who in life have made themselves capable of being helped: ``there is a manner of life which is neither so good as not to require these services after death, nor so bad that such services are of no avail after death.'' From this he concludes: ``Therefore, it is in this life that all the merit or demerit is acquired, which can either relieve or aggravate a man's sufferings after this life.''\endnote{Augustine, \emph{Enchiridion} 110, trans.\ J.~F. Shaw, \emph{Nicene and Post-Nicene Fathers}, first series, vol.~3, New Advent edition, \url{https://www.newadvent.org/fathers/1302.htm} (accessed 16 October 2026).} At his own mother's deathbed he had been asked for this help. Monica, dying at Ostia far from her home, told her sons: ``Lay this body anywhere, let not the care for it trouble you at all. This only I ask, that you will remember me at the Lord's altar, wherever you be.''\endnote{Augustine, \emph{Confessions} IX.11.27, trans.\ J.~G. Pilkington, \emph{Nicene and Post-Nicene Fathers}, first series, vol.~1.}

\subsection{Thirty days of Masses}

Gregory the Great tells, in the fourth book of his \emph{Dialogues}, what happened in his own monastery. A monk named Justus, a physician who had cared for Gregory in his illnesses, was found on his deathbed to have hidden three crowns of gold, against the rule that the monks should own nothing of their own. Gregory, grieved, forbade the brethren to comfort him, so that sorrow might purge his heart before death; and Justus, learning why he had been abandoned, ``straightways sighed for his sin, and in that sorrow gave up the ghost.'' Thirty days after his death Gregory took pity on him:

\begin{quote}
``It is now a good while since that our brother which is departed remaineth in the torments of fire, and therefore we must shew him some charity, and labour what we may to procure his delivery: wherefore go your way, and see that for thirty days following sacrifice be offered for him, so that no one day pass in which, for his absolution and discharge, the healthful sacrifice be not offered.''\endnote{Gregory the Great, \emph{Dialogues} IV.55 (trans.\ P.~W., ed.\ Gardner, 1911), pp.~251--252.}
\end{quote}

On the thirtieth day, Gregory reports, Justus appeared to his brother and said: ``Hitherto have I been in bad case, but now I am well; for this day have I received the communion.'' Gregory concludes that ``the dead monk was by the holy sacrifice delivered from his pains.''\endnote{Gregory, \emph{Dialogues} IV.55 (1911), p.~252.} The apparition is Gregory's report and belongs to pious tradition; what the Church teaches, and the story illustrates, is that the Mass offered for the dead helps them.

\subsection{All Souls}

On the second of November the Church prays at the altar for all the faithful departed. The Council of Trent, which taught that the souls in purgatory are helped above all by the Sacrifice of the altar, charged the bishops to see that ``the suffrages of the faithful who are living, to wit, the sacrifices of masses, prayers, alms, and other works of piety, which have been wont to be performed by the faithful for the other faithful departed, be piously and devoutly performed in accordance with the institutes of the Church.''\endnote{Trent, session~XXV, \emph{Decree concerning Purgatory} (Waterworth), p.~233.}

The Offertory of the Requiem, the same in the 1962 Missal as in the older books, prays for the departed: \latin{Domine Jesu Christe, Rex gloriae, libera animas omnium fidelium defunctorum de poenis inferni, et de profundo lacu: libera eas de ore leonis, ne absorbeat eas tartarus, ne cadant in obscurum: sed signifer sanctus Michael repraesentet eas in lucem sanctam: Quam olim Abrahae promisisti, et semini ejus}.\endnote{\emph{Missale Romanum} (Desclée, 1889), \emph{Missae pro Defunctis}, Offertory, p.~93*.} The Church begs the King of glory to deliver the souls of all the faithful departed from the pains of hell, from the deep pit and the mouth of the lion, that they may not fall into darkness, and that the standard-bearer Saint Michael may bring them into the holy light which God promised of old to Abraham and his seed.

Benedict~XVI writes of this charity: ``The belief that love can reach into the afterlife, that reciprocal giving and receiving is possible, in which our affection for one another continues beyond the limits of death---this has been a fundamental conviction of Christianity throughout the ages and it remains a source of comfort today.'' And he adds: ``It is never too late to touch the heart of another, nor is it ever in vain.''\endnote{\emph{Spe salvi} 48, official English.} In November especially, a Catholic can put this charity into practice: by having Masses offered for his dead, remembering them by name at the altar, giving alms on their behalf, and praying for those whom no one remembers.
